% ============================================================
% CUDA Report - Session 9: Dot Product & Matrix Transpose
% CT225H - Parallel Programming with CUDA
% Nguyễn Thanh Trọng - B2305615
% ============================================================
\documentclass[12pt,a4paper]{article}

% --- Packages ---
\usepackage[utf8]{inputenc}
\usepackage[T5]{fontenc}
\usepackage[vietnamese,english]{babel}
\usepackage[margin=2.5cm]{geometry}
\usepackage{graphicx}
\usepackage{amsmath,amssymb}
\usepackage{booktabs}
\usepackage{listings}
\usepackage{xcolor}
\usepackage{tikz}
\usepackage{pgfplots}
\usepackage{tcolorbox}
\usepackage{fancyhdr}
\usepackage{hyperref}
\usepackage{enumitem}
\usepackage{caption}
\usepackage{float}

\pgfplotsset{compat=1.18}
\usetikzlibrary{arrows.meta, positioning, calc, shapes.geometric, fit, decorations.pathreplacing}

% --- Header/Footer ---
\pagestyle{fancy}
\fancyhf{}
\fancyhead[L]{\small CT225H -- Giữa kỳ}
\fancyhead[R]{\small Nguyễn Thanh Trọng -- B2305615}
\fancyfoot[C]{\thepage}

% --- CUDA Code Style ---
\definecolor{codebg}{RGB}{248,248,248}
\definecolor{codeframe}{RGB}{200,200,200}
\definecolor{keywordcolor}{RGB}{0,0,180}
\definecolor{commentcolor}{RGB}{0,128,0}
\definecolor{stringcolor}{RGB}{163,21,21}
\definecolor{cudakw}{RGB}{180,0,100}

\lstdefinestyle{cudastyle}{
    language=C,
    backgroundcolor=\color{codebg},
    frame=single,
    rulecolor=\color{codeframe},
    basicstyle=\ttfamily\footnotesize,
    keywordstyle=\color{keywordcolor}\bfseries,
    commentstyle=\color{commentcolor}\itshape,
    stringstyle=\color{stringcolor},
    morekeywords={__global__, __shared__, __device__, __host__,
                  blockIdx, blockDim, threadIdx, gridDim,
                  atomicAdd, __syncthreads,
                  cudaMalloc, cudaFree, cudaMemcpy, cudaMemset,
                  cudaMemcpyHostToDevice, cudaMemcpyDeviceToHost,
                  cudaEventCreate, cudaEventRecord, cudaEventSynchronize,
                  cudaEventElapsedTime, cudaEventDestroy,
                  dim3, float},
    keywordstyle=[2]\color{cudakw}\bfseries,
    morekeywords=[2]{__global__, __shared__, __device__, __host__,
                     __syncthreads, atomicAdd},
    numbers=left,
    numberstyle=\tiny\color{gray},
    numbersep=6pt,
    tabsize=4,
    breaklines=true,
    showstringspaces=false,
    captionpos=b,
    xleftmargin=1.5em,
    framexleftmargin=1.5em,
    aboveskip=8pt,
    belowskip=8pt
}

% --- Output Box ---
\newtcolorbox{outputbox}[1][]{
    colback=black!3, colframe=black!50,
    fontupper=\ttfamily\small,
    title={\small\bfseries Output}, #1
}

% --- Title Setup ---
\title{
    \vspace{-1cm}
    \textbf{BÁO CÁO THỰC HÀNH GIỮA KỲ} \\[6pt]
    \Large Lập Trình Song Song Với CUDA \\[4pt]
    \large Dot Product \& Matrix Transpose
}
\author{
    \textbf{Nguyễn Thanh Trọng} -- B2305615 \\
    Môn: CT225H -- Lập trình song song với CUDA
}
\date{\today}

% ============================================================
\begin{document}
\maketitle
\thispagestyle{fancy}

% ============================================================
\section{Giới thiệu}
% ============================================================

Báo cáo trình bày việc triển khai hai bài toán cơ bản trong tính toán song song trên GPU bằng CUDA:

\begin{enumerate}[leftmargin=2em]
    \item \textbf{Dot Product} (tích vô hướng): Tính $\displaystyle \text{result} = \sum_{i=0}^{N-1} A_i \times B_i$ với $N = 1{,}000{,}000$ phần tử. Sử dụng kỹ thuật \textbf{Parallel Reduction} kết hợp Shared Memory.
    \item \textbf{Matrix Transpose} (chuyển vị ma trận): Tính $A^T$ với $A^T[j][i] = A[i][j]$, kích thước $1024 \times 1024$. So sánh phương pháp Naive và Shared Memory Tiling.
\end{enumerate}

\noindent\textbf{Môi trường thực thi:} Google Colab -- GPU Tesla T4 (Turing), CUDA 13.0, Driver 580.82.07.

% ============================================================
\section{Bài toán 1: Dot Product với Parallel Reduction}
% ============================================================

\subsection{Phát biểu bài toán}

Cho hai vector $A = [a_0, a_1, \ldots, a_{N-1}]$ và $B = [b_0, b_1, \ldots, b_{N-1}]$ với $N = 10^6$. Tính:
\[
    \text{result} = A \cdot B = \sum_{i=0}^{N-1} a_i \times b_i
\]

Đây là phép toán \textit{reduction} -- rút gọn $N$ phần tử thành 1 giá trị duy nhất. Trên CPU thao tác này là tuần tự $O(N)$, nhưng trên GPU có thể song song hóa xuống $O(\log N)$ theo chiều sâu.

\subsection{Thiết kế thuật toán CUDA}

\subsubsection{Mô hình tổng quan}

\begin{figure}[H]
\centering
\begin{tikzpicture}[
    node distance=0.6cm and 0.4cm,
    block/.style={draw, rounded corners, fill=blue!10, minimum width=2.4cm, minimum height=0.7cm, align=center, font=\small},
    arrow/.style={-{Stealth[length=5pt]}, thick}
]
    % Phase 1
    \node[block, fill=orange!15] (mul) {Nhân từng cặp\\$c_i = a_i \times b_i$};
    \node[block, fill=green!12, below=of mul] (grid) {Grid-stride loop\\gom nhiều $c_i$ / thread};
    \node[block, fill=cyan!12, below=of grid] (shared) {Ghi vào\\Shared Memory};
    \node[block, fill=yellow!15, below=of shared] (reduce) {Tree-based\\Reduction};
    \node[block, fill=red!12, below=of reduce] (atomic) {atomicAdd\\kết quả block $\to$ global};

    \draw[arrow] (mul) -- (grid);
    \draw[arrow] (grid) -- (shared);
    \draw[arrow] (shared) -- (reduce);
    \draw[arrow] (reduce) -- (atomic);

    % Labels
    \node[left=0.8cm of mul, font=\footnotesize\itshape, text=gray] {Phase 1};
    \node[left=0.8cm of shared, font=\footnotesize\itshape, text=gray] {Phase 2};
    \node[left=0.8cm of atomic, font=\footnotesize\itshape, text=gray] {Phase 3};
\end{tikzpicture}
\caption{Pipeline xử lý Dot Product trên GPU}
\end{figure}

\subsubsection{Chi tiết từng giai đoạn}

\paragraph{Phase 1 -- Nhân phần tử \& Grid-stride loop.}
Mỗi thread được gán một chỉ số khởi đầu \texttt{idx = blockIdx.x * blockDim.x + threadIdx.x}. Thread duyệt qua nhiều phần tử bằng \textbf{grid-stride loop}: tăng \texttt{idx} thêm \texttt{blockDim.x * gridDim.x} mỗi lần lặp, tích lũy $c_i = a_i \times b_i$ vào biến cục bộ \texttt{temp}. Kỹ thuật này đảm bảo kernel hoạt động đúng với mọi kích thước dữ liệu, không phụ thuộc số block.

\paragraph{Phase 2 -- Shared Memory Reduction.}
Kết quả \texttt{temp} được ghi vào mảng \texttt{\_\_shared\_\_ cache[BLOCK\_SIZE]}. Sau \texttt{\_\_syncthreads()}, thực hiện \textbf{tree-based reduction}: ở mỗi bước, stride giảm một nửa ($128 \to 64 \to \ldots \to 1$), thread \texttt{tid < stride} cộng \texttt{cache[tid] += cache[tid + stride]}.

\begin{figure}[H]
\centering
\begin{tikzpicture}[
    cell/.style={draw, minimum width=0.7cm, minimum height=0.55cm, font=\scriptsize},
    highlight/.style={cell, fill=blue!15},
    arrow/.style={-{Stealth[length=3pt]}, gray, thin}
]
    % Level 0 - 8 elements
    \node[font=\footnotesize\bfseries] at (-1.5, 0) {Step 0};
    \foreach \i/\v in {0/a,1/b,2/c,3/d,4/e,5/f,6/g,7/h} {
        \node[cell] (L0-\i) at (\i*0.8, 0) {\v};
    }
    
    % Level 1 - stride=4
    \node[font=\footnotesize\bfseries] at (-1.5, -1.2) {Step 1};
    \foreach \i/\v in {0/{a+e},1/{b+f},2/{c+g},3/{d+h}} {
        \node[highlight] (L1-\i) at (\i*0.8, -1.2) {\v};
    }
    \foreach \i in {4,5,6,7} {
        \node[cell, fill=gray!10] (L1-\i) at (\i*0.8, -1.2) {--};
    }
    \foreach \i [evaluate=\i as \j using int(\i+4)] in {0,1,2,3} {
        \draw[arrow] (L0-\i.south) -- (L1-\i.north);
        \draw[arrow] (L0-\j.south) -- (L1-\i.north);
    }
    
    % Level 2 - stride=2
    \node[font=\footnotesize\bfseries] at (-1.5, -2.4) {Step 2};
    \foreach \i/\v in {0/{$\Sigma_{0..5}$},1/{$\Sigma_{1..6}$}} {
        \node[highlight, fill=green!15] (L2-\i) at (\i*0.8, -2.4) {\v};
    }
    \foreach \i in {2,3,4,5,6,7} {
        \node[cell, fill=gray!10] (L2-\i) at (\i*0.8, -2.4) {--};
    }
    \draw[arrow] (L1-0.south) -- (L2-0.north);
    \draw[arrow] (L1-2.south) -- (L2-0.north);
    \draw[arrow] (L1-1.south) -- (L2-1.north);
    \draw[arrow] (L1-3.south) -- (L2-1.north);

    % Level 3 - stride=1
    \node[font=\footnotesize\bfseries] at (-1.5, -3.6) {Step 3};
    \node[highlight, fill=red!15] (L3-0) at (0*0.8, -3.6) {$\Sigma$};
    \foreach \i in {1,2,3,4,5,6,7} {
        \node[cell, fill=gray!10] (L3-\i) at (\i*0.8, -3.6) {--};
    }
    \draw[arrow] (L2-0.south) -- (L3-0.north);
    \draw[arrow] (L2-1.south) -- (L3-0.north);
\end{tikzpicture}
\caption{Tree-based parallel reduction (minh họa 8 phần tử, BLOCK\_SIZE = 256 tương tự với 8 bước)}
\end{figure}

\paragraph{Phase 3 -- atomicAdd.}
Thread 0 của mỗi block dùng \texttt{atomicAdd(result, cache[0])} để cộng kết quả block vào biến global. Với $N = 10^6$ và BLOCK\_SIZE = 256, có $\lceil 10^6 / 256 \rceil = 3907$ block, tức 3907 lần atomicAdd -- chấp nhận được về mặt hiệu năng.

\subsection{Mã nguồn CUDA}

\begin{lstlisting}[style=cudastyle, caption={Kernel Dot Product -- Parallel Reduction}]
#define N 1000000
#define BLOCK_SIZE 256

__global__ void dotProductKernel(float *a, float *b, 
                                  float *result, int n) {
    __shared__ float cache[BLOCK_SIZE];
    int tid = threadIdx.x;
    int idx = blockIdx.x * blockDim.x + threadIdx.x;

    // Phase 1: Grid-stride loop -- moi thread tich luy nhieu cap
    float temp = 0.0f;
    while (idx < n) {
        temp += a[idx] * b[idx];
        idx += blockDim.x * gridDim.x;
    }
    cache[tid] = temp;
    __syncthreads();

    // Phase 2: Tree-based reduction trong shared memory
    for (int stride = blockDim.x / 2; stride > 0; stride >>= 1) {
        if (tid < stride)
            cache[tid] += cache[tid + stride];
        __syncthreads();
    }

    // Phase 3: Block result -> global
    if (tid == 0)
        atomicAdd(result, cache[0]);
}
\end{lstlisting}

\begin{lstlisting}[style=cudastyle, caption={Host code -- Dot Product}, firstnumber=30]
int numBlocks = (N + BLOCK_SIZE - 1) / BLOCK_SIZE;
cudaMemset(d_result, 0, sizeof(float));

cudaEvent_t start, stop;
cudaEventCreate(&start);
cudaEventCreate(&stop);

cudaEventRecord(start);
dotProductKernel<<<numBlocks, BLOCK_SIZE>>>(d_a, d_b, d_result, N);
cudaEventRecord(stop);
cudaEventSynchronize(stop);

float gpu_time;
cudaEventElapsedTime(&gpu_time, start, stop);
\end{lstlisting}

\subsection{Kết quả thực thi}

\begin{outputbox}
=== Dot Product (N = 1000000) ===
[CPU] Result: 244166.156250
[CPU] Time:   3.0210 ms
[GPU] Result: 244218.968750
[GPU] Time:   0.1928 ms
Speedup: 15.67x
Absolute error: 52.812500
\end{outputbox}

\subsection{Phân tích kết quả}

\begin{table}[H]
\centering
\begin{tabular}{lrrr}
\toprule
\textbf{Phương pháp} & \textbf{Thời gian (ms)} & \textbf{Speedup} & \textbf{Kết quả} \\
\midrule
CPU (tuần tự)   & 3.0210 & 1.00$\times$ & 244166.16 \\
GPU (Reduction) & 0.1928 & 15.67$\times$ & 244218.97 \\
\bottomrule
\end{tabular}
\caption{So sánh hiệu năng Dot Product}
\end{table}

\begin{figure}[H]
\centering
\begin{tikzpicture}
\begin{axis}[
    ybar,
    width=8cm, height=5cm,
    ylabel={Thời gian (ms)},
    symbolic x coords={CPU, GPU},
    xtick=data,
    ymin=0, ymax=4,
    bar width=1.5cm,
    nodes near coords,
    nodes near coords align={vertical},
    every node near coord/.append style={font=\small},
    enlarge x limits=0.5,
]
\addplot[fill=red!40]   coordinates {(CPU, 3.0210)};
\addplot[fill=green!50] coordinates {(GPU, 0.1928)};
\legend{CPU, GPU}
\end{axis}
\end{tikzpicture}
\caption{So sánh thời gian CPU vs GPU -- Dot Product}
\end{figure}

\paragraph{Sai số floating-point.}
Sai số tuyệt đối $\approx 52.81$ (sai số tương đối $\approx 0.02\%$) là do thứ tự cộng dồn khác nhau giữa CPU (tuần tự) và GPU (parallel reduction). Phép cộng \texttt{float} không có tính kết hợp (non-associative), nên kết quả phụ thuộc thứ tự tính toán. Đây là hiện tượng bình thường, không phải lỗi thuật toán.

% ============================================================
\section{Bài toán 2: Matrix Transpose}
% ============================================================

\subsection{Phát biểu bài toán}

Cho ma trận $A$ kích thước $N \times M$ ($N = M = 1024$), lưu dạng row-major trong bộ nhớ liên tục. Tính ma trận chuyển vị:
\[
    A^T[j][i] = A[i][j], \quad \forall\, 0 \le i < N,\; 0 \le j < M
\]

\subsection{Thiết kế thuật toán}

Hai phiên bản GPU được triển khai để so sánh ảnh hưởng của memory access pattern lên hiệu năng.

\subsubsection{Naive Transpose}

\begin{figure}[H]
\centering
\begin{tikzpicture}[
    cell/.style={draw, minimum width=0.7cm, minimum height=0.7cm, font=\scriptsize},
    arrow/.style={-{Stealth[length=5pt]}, thick}
]
    % Input matrix
    \node[font=\small\bfseries] at (1.4, 1) {Input A (row-major)};
    \foreach \i in {0,...,3} {
        \foreach \j in {0,...,3} {
            \pgfmathtruncatemacro{\idx}{\i*4+\j}
            \node[cell, fill=blue!10] (A\i\j) at (\j*0.8, -\i*0.8) {\idx};
        }
    }
    
    % Arrow
    \draw[arrow, very thick, blue!60] (3.8, -1.2) -- (5.2, -1.2) 
        node[midway, above, font=\scriptsize] {Transpose};
    
    % Output matrix
    \node[font=\small\bfseries] at (7.6, 1) {Output $A^T$};
    \foreach \i in {0,...,3} {
        \foreach \j in {0,...,3} {
            \pgfmathtruncatemacro{\idx}{\j*4+\i}
            \node[cell, fill=orange!15] (B\i\j) at (\j*0.8+6, -\i*0.8) {\idx};
        }
    }
    
    % Highlight coalesced/non-coalesced
    \draw[red, thick, dashed] (-0.05, 0.35) rectangle (2.85, -0.35);
    \node[font=\tiny, red, right] at (3, 0) {Read: coalesced};
    
    \draw[red, thick, dashed] (5.95, 0.35) rectangle (6.55, -2.75);
    \node[font=\tiny, red, below] at (6.25, -3) {Write: strided!};
\end{tikzpicture}
\caption{Naive Transpose -- đọc coalesced theo hàng, ghi strided theo cột}
\end{figure}

Mỗi thread đọc \texttt{in[(y+j)*cols + x]} (liên tiếp trong bộ nhớ -- \textbf{coalesced read}) và ghi vào \texttt{out[x*rows + (y+j)]} (bước nhảy lớn -- \textbf{non-coalesced write}). Non-coalesced write gây nhiều transaction trên global memory, giảm throughput đáng kể.

\subsubsection{Shared Memory Transpose (Optimized)}

\begin{figure}[H]
\centering
\begin{tikzpicture}[
    box/.style={draw, rounded corners, minimum width=3cm, minimum height=1cm, align=center, font=\small},
    arrow/.style={-{Stealth[length=5pt]}, thick}
]
    \node[box, fill=blue!10] (read) at (0, 0) {Coalesced Read\\Global $\to$ Shared};
    \node[box, fill=yellow!15] (tile) at (5, 0) {Shared Memory\\Tile [32][33]};
    \node[box, fill=green!12] (write) at (10, 0) {Coalesced Write\\Shared $\to$ Global};
    
    \draw[arrow] (read) -- (tile) node[midway, above, font=\scriptsize] {Row-wise};
    \draw[arrow] (tile) -- (write) node[midway, above, font=\scriptsize] {Column-wise};
    
    \node[font=\scriptsize, text=gray, below=0.2cm of tile] {Padding +1 tránh bank conflict};
\end{tikzpicture}
\caption{Shared Memory Transpose -- cả read và write đều coalesced}
\end{figure}

Ý tưởng chính:
\begin{enumerate}[leftmargin=2em]
    \item \textbf{Coalesced Read:} Một tile $32 \times 32$ được đọc từ global memory vào shared memory theo hàng (coalesced).
    \item \textbf{Hoán đổi chỉ số trong shared memory:} Đọc lại từ shared memory theo cột (truy cập nhanh, ~1 cycle latency).
    \item \textbf{Coalesced Write:} Ghi ra vị trí mới trong global memory theo hàng (coalesced).
\end{enumerate}

\paragraph{Bank conflict \& Padding.}
Shared memory được chia thành 32 bank. Khi 32 thread cùng truy cập các phần tử cách nhau đúng 32 vị trí (cùng cột), xảy ra \textbf{bank conflict} -- truy cập bị tuần tự hóa. Giải pháp: khai báo \texttt{tile[32][33]} thay vì \texttt{tile[32][32]}. Phần tử padding (+1) làm lệch chỉ số, đảm bảo mỗi thread truy cập bank khác nhau.

\subsection{Mã nguồn CUDA}

\begin{lstlisting}[style=cudastyle, caption={Kernel Naive Transpose}]
#define TILE_DIM 32
#define BLOCK_ROWS 8

__global__ void transposeNaive(float *in, float *out, 
                                int rows, int cols) {
    int x = blockIdx.x * TILE_DIM + threadIdx.x;
    int y = blockIdx.y * TILE_DIM + threadIdx.y;

    for (int j = 0; j < TILE_DIM; j += BLOCK_ROWS) {
        if (x < cols && (y + j) < rows)
            out[x * rows + (y + j)] = in[(y + j) * cols + x];
    }
}
\end{lstlisting}

\begin{lstlisting}[style=cudastyle, caption={Kernel Shared Memory Transpose}]
__global__ void transposeSharedMem(float *in, float *out, 
                                    int rows, int cols) {
    __shared__ float tile[TILE_DIM][TILE_DIM + 1]; // +1 padding

    int x = blockIdx.x * TILE_DIM + threadIdx.x;
    int y = blockIdx.y * TILE_DIM + threadIdx.y;

    // Coalesced read: global -> shared
    for (int j = 0; j < TILE_DIM; j += BLOCK_ROWS) {
        if (x < cols && (y + j) < rows)
            tile[threadIdx.y + j][threadIdx.x] = in[(y+j)*cols + x];
    }
    __syncthreads();

    // Doi block index de ghi ra vi tri chuyen vi
    x = blockIdx.y * TILE_DIM + threadIdx.x;
    y = blockIdx.x * TILE_DIM + threadIdx.y;

    // Coalesced write: shared -> global
    for (int j = 0; j < TILE_DIM; j += BLOCK_ROWS) {
        if (x < rows && (y + j) < cols)
            out[(y+j)*rows + x] = tile[threadIdx.x][threadIdx.y + j];
    }
}
\end{lstlisting}

\begin{lstlisting}[style=cudastyle, caption={Host code -- cấu hình grid \& block}, firstnumber=25]
dim3 block(TILE_DIM, BLOCK_ROWS);   // (32, 8) = 256 threads/block
dim3 grid((cols + TILE_DIM - 1) / TILE_DIM,
          (rows + TILE_DIM - 1) / TILE_DIM);  // (32, 32) blocks

// Naive
transposeNaive<<<grid, block>>>(d_input, d_output, rows, cols);

// Shared Memory
transposeSharedMem<<<grid, block>>>(d_input, d_output, rows, cols);
\end{lstlisting}

\subsection{Kết quả thực thi}

\begin{outputbox}
=== Matrix Transpose (1024 x 1024) ===
[CPU]          Time: 19.7390 ms
[GPU Naive]    Time: 0.2270 ms  |  Correct: Yes
[GPU Shared]   Time: 0.0619 ms  |  Correct: Yes

--- Performance Summary ---
CPU Time:              19.7390 ms
GPU Naive Time:        0.2270 ms  (Speedup: 86.94x vs CPU)
GPU Shared Mem Time:   0.0619 ms  (Speedup: 318.95x vs CPU)
Shared vs Naive:       3.67x faster
\end{outputbox}

\subsection{Phân tích kết quả}

\begin{table}[H]
\centering
\begin{tabular}{lrrl}
\toprule
\textbf{Phương pháp} & \textbf{Thời gian (ms)} & \textbf{Speedup vs CPU} & \textbf{Đúng?} \\
\midrule
CPU (tuần tự)         & 19.7390 & 1.00$\times$   & -- \\
GPU Naive             &  0.2270 & 86.94$\times$  & Yes \\
GPU Shared Memory     &  0.0619 & 318.95$\times$ & Yes \\
\bottomrule
\end{tabular}
\caption{So sánh hiệu năng Matrix Transpose ($1024 \times 1024$)}
\end{table}

\begin{figure}[H]
\centering
\begin{tikzpicture}
\begin{axis}[
    ybar,
    width=10cm, height=5.5cm,
    ylabel={Thời gian (ms)},
    symbolic x coords={CPU, GPU Naive, GPU Shared},
    xtick=data,
    ymin=0,
    bar width=1.2cm,
    nodes near coords,
    every node near coord/.append style={font=\small},
    enlarge x limits=0.35,
    legend style={at={(0.98,0.95)}, anchor=north east},
]
\addplot[fill=red!40]    coordinates {(CPU, 19.739)};
\addplot[fill=orange!50] coordinates {(GPU Naive, 0.227)};
\addplot[fill=green!50]  coordinates {(GPU Shared, 0.0619)};
\end{axis}
\end{tikzpicture}
\caption{So sánh thời gian -- Matrix Transpose}
\end{figure}

\begin{figure}[H]
\centering
\begin{tikzpicture}
\begin{axis}[
    ybar,
    width=8cm, height=5cm,
    ylabel={Speedup so với CPU ($\times$)},
    symbolic x coords={GPU Naive, GPU Shared},
    xtick=data,
    ymin=0, ymax=400,
    bar width=1.5cm,
    nodes near coords,
    every node near coord/.append style={font=\small\bfseries},
    enlarge x limits=0.5,
]
\addplot[fill=orange!50] coordinates {(GPU Naive, 86.94)};
\addplot[fill=green!50]  coordinates {(GPU Shared, 318.95)};
\end{axis}
\end{tikzpicture}
\caption{Speedup so với CPU}
\end{figure}

\paragraph{Tại sao Shared Memory nhanh hơn Naive 3.67$\times$?}

\begin{itemize}[leftmargin=2em]
    \item \textbf{Naive:} Ghi non-coalesced -- mỗi warp (32 thread) ghi vào 32 địa chỉ cách nhau $1024 \times 4 = 4096$ byte. GPU phải thực hiện tối đa 32 transaction riêng biệt (mỗi transaction 32 byte) thay vì 1 transaction 128 byte.
    \item \textbf{Shared Memory:} Cả đọc và ghi global memory đều coalesced. Shared memory có latency $\sim$1--2 cycle (vs $\sim$200--400 cycle cho global memory cache miss). Padding \texttt{+1} loại bỏ bank conflict hoàn toàn.
    \item \textbf{Tổng hợp:} Shared Memory tận dụng tối đa memory bandwidth của GPU, đạt speedup 318.95$\times$ so với CPU.
\end{itemize}

% ============================================================
\section{Tổng hợp \& So sánh}
% ============================================================

\begin{table}[H]
\centering
\begin{tabular}{llrrr}
\toprule
\textbf{Bài toán} & \textbf{Phương pháp} & \textbf{CPU (ms)} & \textbf{GPU (ms)} & \textbf{Speedup} \\
\midrule
\multirow{1}{*}{Dot Product}     & Parallel Reduction & 3.02  & 0.19  & 15.67$\times$ \\
\midrule
\multirow{2}{*}{Matrix Transpose} & Naive              & 19.74 & 0.23  & 86.94$\times$ \\
                                   & Shared Memory      & 19.74 & 0.062 & 318.95$\times$ \\
\bottomrule
\end{tabular}
\caption{Tổng hợp hiệu năng tất cả bài toán}
\end{table}

\paragraph{Nhận xét:}
\begin{itemize}[leftmargin=2em]
    \item Matrix Transpose đạt speedup cao hơn Dot Product vì bản chất bài toán ``embarrassingly parallel'' -- mỗi phần tử xử lý độc lập, không cần reduction.
    \item Dot Product bị giới hạn bởi bước atomicAdd cuối cùng (serialization point), nhưng vẫn đạt 15.67$\times$ nhờ tận dụng parallel reduction.
    \item Shared memory là kỹ thuật tối ưu thiết yếu trên GPU -- cải thiện 3.67$\times$ so với naive chỉ bằng cách thay đổi memory access pattern.
\end{itemize}

% ============================================================
\section{Kết luận}
% ============================================================

Bài thực hành đã triển khai thành công hai thuật toán song song trên GPU CUDA:

\begin{enumerate}[leftmargin=2em]
    \item \textbf{Dot Product:} Áp dụng parallel reduction với shared memory và atomicAdd, đạt speedup 15.67$\times$ so với CPU.
    \item \textbf{Matrix Transpose:} So sánh naive vs shared memory tiling, chứng minh ảnh hưởng lớn của coalesced memory access (speedup từ 86.94$\times$ lên 318.95$\times$).
\end{enumerate}

Các kỹ thuật tối ưu chính được áp dụng: \textbf{grid-stride loop}, \textbf{shared memory}, \textbf{tree-based reduction}, \textbf{tile-based processing}, và \textbf{bank conflict avoidance (padding +1)}.

\end{document}
